% GENERATED FILE. Edit canonical lessons, then run npm run generate:books.
% canonical-source-hash: fnv1a64:c4d6df4d1e38a8aa
% canonical-lessons: PT-C496-ternura, PT-C496-carinhoso, PT-C496-afeto, PT-C496-acariciar, PT-C496-mimar

\chapter{Ternura e carinho}
\label{ch:pt-ternura-e-carinho}
\hlchaptermodality{496}

\begin{chapteropening}
\textbf{By the end of this chapter:} \emph{I can talk about tenderness and affection.}

Complete the last lesson of chapter 496: I can talk about tenderness and affection.
\end{chapteropening}

\section[a ternura]{a ternura — tenderness}

\label{lesson:PT-C496-ternura}



\begin{quote}
Before the new one: say the Portuguese for (monthly) pocket money, then the Portuguese for spoiled.
\end{quote}

\subsection*{You'll want to know: a ternura}
\textbf{a ternura} — \textquotedblleft{}tenderness\textquotedblright{}.

\begin{quote}
A avó olhava para o neto com ternura.
\end{quote}

\begin{quote}
\emph{Grandma was looking at her grandson with tenderness.}
\end{quote}

Built on \emph{terno} \textquotedblleft{}tender\textquotedblright{}.

\begin{cousinweb}
\emph{ternura} is built on \emph{terno} \textquotedblleft{}tender\textquotedblright{}, which is inherited from Latin \emph{tener} \textquotedblleft{}soft, tender\textquotedblright{} (Lewis \& Short).
\end{cousinweb}

\begin{grammarlens}[title={What you've built}]
One more word for telling, explaining and describing at B1.
\end{grammarlens}

\subsection*{Your turn}
\begin{itemize}
\raggedright
  \item \emph{Say it:} \emph{a ternura}
  \item \emph{Say it:} \emph{a ternura}, once more
  \item \emph{Say it:} a sentence of your own with \emph{a ternura}
  \item \emph{Recall:} say the Portuguese for (monthly) pocket money, then the Portuguese for spoiled, then say \emph{a ternura} again
\end{itemize}

\begin{tcolorbox}[breakable,colback=teal!4,colframe=teal!35!black,title={Before you move on}]
What does a ternura mean? (\textquotedblleft{}Tenderness\textquotedblright{}.) Say it once more.
\end{tcolorbox}
\section[carinhoso]{carinhoso — affectionate, loving}

\label{lesson:PT-C496-carinhoso}



\begin{quote}
Before the new one: say the Portuguese for spoiled, then the Portuguese for tenderness.
\end{quote}

\subsection*{You'll want to know: carinhoso}
\textbf{carinhoso} — \textquotedblleft{}affectionate, loving\textquotedblright{}.

\begin{quote}
O meu filho é muito carinhoso.
\end{quote}

\begin{quote}
\emph{My son is very affectionate.}
\end{quote}

Built on \emph{o carinho} — affection, a caress.

\begin{cousinweb}
\emph{carinhoso} is built on \emph{carinho} \textquotedblleft{}affection\textquotedblright{}, which is usually traced to Spanish \emph{cariño} (EtymDB); the deeper origin is not settled.
\end{cousinweb}

\begin{grammarlens}[title={What you've built}]
One more word for telling, explaining and describing at B1.
\end{grammarlens}

\subsection*{Your turn}
\begin{itemize}
\raggedright
  \item \emph{Say it:} \emph{carinhoso}
  \item \emph{Say it:} \emph{carinhoso}, once more
  \item \emph{Say it:} a sentence of your own with \emph{carinhoso}
  \item \emph{Recall:} say the Portuguese for spoiled, then the Portuguese for tenderness, then say \emph{carinhoso} again
\end{itemize}

\begin{tcolorbox}[breakable,colback=teal!4,colframe=teal!35!black,title={Before you move on}]
What does carinhoso mean? (\textquotedblleft{}Affectionate, loving\textquotedblright{}.) Say it once more.
\end{tcolorbox}
\section[o afeto]{o afeto — affection}

\label{lesson:PT-C496-afeto}



\begin{quote}
Before the new one: say the Portuguese for tenderness, then the Portuguese for affectionate.
\end{quote}

\subsection*{You'll want to know: o afeto}
\textbf{o afeto} — \textquotedblleft{}affection\textquotedblright{}.

\begin{quote}
Os animais também precisam de afeto.
\end{quote}

\begin{quote}
\emph{Animals need affection too.}
\end{quote}

Before AO90 spelled \emph{afecto} in Portugal.

\begin{cousinweb}
\emph{afeto} comes from Latin \emph{affectus} \textquotedblleft{}a state of mind, a feeling\textquotedblright{} (EtymDB), from \emph{afficere} \textquotedblleft{}to act on\textquotedblright{}.
\end{cousinweb}

\begin{grammarlens}[title={What you've built}]
One more word for telling, explaining and describing at B1.
\end{grammarlens}

\subsection*{Your turn}
\begin{itemize}
\raggedright
  \item \emph{Say it:} \emph{o afeto}
  \item \emph{Say it:} \emph{o afeto}, once more
  \item \emph{Say it:} a sentence of your own with \emph{o afeto}
  \item \emph{Recall:} say the Portuguese for tenderness, then the Portuguese for affectionate, then say \emph{o afeto} again
\end{itemize}

\begin{tcolorbox}[breakable,colback=teal!4,colframe=teal!35!black,title={Before you move on}]
What does o afeto mean? (\textquotedblleft{}Affection\textquotedblright{}.) Say it once more.
\end{tcolorbox}
\section[acariciar]{acariciar — to stroke, to caress}

\label{lesson:PT-C496-acariciar}



\begin{quote}
Before the new one: say the Portuguese for affectionate, then the Portuguese for affection.
\end{quote}

\subsection*{You'll want to know: acariciar}
\textbf{acariciar} — \textquotedblleft{}to stroke, to caress\textquotedblright{}.

\begin{quote}
A menina acariciava o gato.
\end{quote}

\begin{quote}
\emph{The girl was stroking the cat.}
\end{quote}

A caress is \emph{uma carícia}.

\begin{cousinweb}
\emph{acariciar} is built on \emph{carícia} \textquotedblleft{}a caress\textquotedblright{}, from Italian \emph{carezza}, which goes back to Latin \emph{cārus} \textquotedblleft{}dear\textquotedblright{} (Lewis \& Short).
\end{cousinweb}

\begin{grammarlens}[title={What you've built}]
One more word for telling, explaining and describing at B1.
\end{grammarlens}

\subsection*{Your turn}
\begin{itemize}
\raggedright
  \item \emph{Say it:} \emph{acariciar}
  \item \emph{Say it:} \emph{acariciar}, once more
  \item \emph{Say it:} a sentence of your own with \emph{acariciar}
  \item \emph{Recall:} say the Portuguese for affectionate, then the Portuguese for affection, then say \emph{acariciar} again
\end{itemize}

\begin{tcolorbox}[breakable,colback=teal!4,colframe=teal!35!black,title={Before you move on}]
What does acariciar mean? (\textquotedblleft{}To stroke, to caress\textquotedblright{}.) Say it once more.
\end{tcolorbox}
\section[mimar]{mimar — to spoil, to pamper}

\label{lesson:PT-C496-mimar}



\begin{quote}
Before the new one: say the Portuguese for affection, then the Portuguese for to stroke.
\end{quote}

\subsection*{You'll want to know: mimar}
\textbf{mimar} — \textquotedblleft{}to spoil, to pamper\textquotedblright{}.

\begin{quote}
Os avós mimam os netos com doces.
\end{quote}

\begin{quote}
\emph{The grandparents spoil their grandchildren with sweets.}
\end{quote}

\emph{Mimado} — spoiled.

\begin{cousinweb}
\emph{mimar} is built on \emph{mimo} \textquotedblleft{}a treat, a caress\textquotedblright{}, whose origin is not settled.
\end{cousinweb}

\begin{grammarlens}[title={What you've built}]
One more word for telling, explaining and describing at B1.
\end{grammarlens}

\subsection*{Your turn}
\begin{itemize}
\raggedright
  \item \emph{Say it:} \emph{mimar}
  \item \emph{Say it:} \emph{mimar}, once more
  \item \emph{Say it:} a sentence of your own with \emph{mimar}
  \item \emph{Recall:} say the Portuguese for affection, then the Portuguese for to stroke, then say \emph{mimar} again
\end{itemize}

\begin{tcolorbox}[breakable,colback=teal!4,colframe=teal!35!black,title={Before you move on}]
What does mimar mean? (\textquotedblleft{}To spoil, to pamper\textquotedblright{}.) Say it once more.
\end{tcolorbox}
